\section{\ref{sec:acet}장. \nt{ACET}를 더하다}

아세틸콜린의 세 가지 작용은 절편과 생체 뇌에서 측정되었고, 쓰기와 읽기의 충돌은 이 책의 모델에서 보였다. \nt{ACET}가 쓰기 허용 선으로 쓰이며 예상된 불확실성을 알린다는 것은 실험의 부분적 뒷받침이 있는 가설이다.

{\footnotesize
\setlength{\tabcolsep}{3pt}\renewcommand{\arraystretch}{1.15}
\begin{longtable}{>{\raggedright}p{0.5cm}>{\raggedright}p{4.0cm}>{\raggedright}p{5.6cm}>{\raggedright}p{2.3cm}>{\raggedright\arraybackslash}p{1.7cm}}
\toprule
No. & 명제 & 실험과 측정 내용 & 출처 & 상태 \\
\midrule\endfirsthead
\toprule
No. & 명제 & 실험과 측정 내용 & 출처 & 상태 \\
\midrule\endhead
1 & 뇌의 \nt{ACET} 뉴런은 많지 않고 몇 개의 집단에 모여 있으며, 출력은 피질 전체로 퍼진다. 브로드캐스트 채널이다 & 쥐에서 아세틸콜린 합성 효소에 대한 항체 염색과 역행성 표지: 피질, 해마, 편도체, 후각 망울, 시상은 아세틸콜린을 주로 기저 전뇌와 뇌간의 조밀한 세포 집단 여섯 개에서 받는다 & \cite{Mesulam1983} & 측정됨 \\
2 & 피질의 아세틸콜린 수준은 깨어 있을 때 높고 서파 수면에서 낮다 & 자유롭게 움직이는 고양이의 피질과 해마에서 뇌파 기록과 함께 한 미세투석: 아세틸콜린 방출은 서파 수면보다 조용히 깨어 있을 때 $\sim100\%$, 활발히 깨어 있을 때 $\sim175\%$ 높다. 렘수면에서 피질의 방출은 깨어 있을 때와 같다 & \cite{Marrosu1995,Hasselmo1999} & 측정됨 \\
3 & 신호의 의미는 수용체가 정한다. 아세틸콜린에는 빠른 채널형 수용체와 세포 안 반응을 일으키는 느린 수용체가 있다(명제~\ref{prop:receptor}) & 측정 결과의 총설: 아세틸콜린이 흥분성, 전달, 가소성에 미치는 작용은 방출 장소, 수용체 아형(니코틴성은 이온 채널, 무스카린성은 G 단백질 경유), 수신 세포에 달려 있다 & \cite{Picciotto2012} & 측정됨 \\
4 & 첫째 작용: 외부 입력은 거의 건드리지 않고 내부 연결을 약하게 한다(식~\eqref{eq:acetnet}의 인수 $1-a$) & 쥐 후각 피질 절편: 아세틸콜린 작용제 $100$\,\textmu M에서 같은 세포의 내부 연결(Ib층) 전위는 $60\%$ 넘게, 외부 입력(Ia층) 전위는 $15\%$ 미만으로 줄어든다. 억제는 시냅스 전 작용이다. 해마 절편(CA1): 내부 경로 $89\%$ 대 입력 경로 $40\%$ & \cite{HasselmoBower1992,HasselmoSchnell1994} & 측정됨 \\
5 & 둘째 작용: 입력을 강하게 한다(인수 $\frac{1+a}{2}$) & 신피질 절편: 니코틴성 수용체는 시상에서 오는 입력 시냅스를 강화하고 피질 내부 시냅스에는 작용하지 않는다. 무스카린성 수용체는 두 종류를 모두 약하게 한다. 아세틸콜린은 뉴런의 흥분성을 높인다(총설) & \cite{Gil1997,Hasselmo2006} & 측정됨 \\
6 & 셋째 작용: 가중치 변화를 쉽게 한다(학습률 $\eta a$) & 쥐: 피질 아세틸콜린의 공급원인 기저핵을 전기 자극하면서 소리를 함께 들려주면, 다 자란 동물에서도 청각 피질 지도가 들려준 소리 쪽으로 크게 재편된다 & \cite{KilgardMerzenich1998,Hasselmo2006} & 측정됨 \\
7 & 식~\eqref{eq:acetnet}의 두 번째 식, 즉 드문 패턴용 헤브 규칙은 일부로부터 패턴을 완성하는 연상 메모리를 준다 & 성긴 패턴과 공분산 규칙을 쓰는 네트워크의 용량 계산: 활성 뉴런 비율 $p$가 작으면 네트워크는 $p=1/2$일 때보다 눈에 띄게 많은 패턴을 저장한다 & \cite{Tsodyks1988} & 이론 \\
8 & 낮은 $a$에서의 쓰기는 메모리를 망가뜨린다. 네트워크는 본 것이 아니라 떠올린 것을 기록한다. $a=0.1$에서도 손상은 $a=0$보다 약하지 않다 & \texttt{models/\allowbreak acet\_memory.py}: 뉴런 $200$개, $20$개짜리 패턴 다섯 개, 새 패턴은 그중 하나와 뉴런 $10$개를 공유, 실행 $10$회. $a=0$에서 새 패턴은 기억되지 않고, 옛 패턴은 남의 뉴런 $10$개 중 $5.9$개를 끌고 나온다. $a=0.1$에서 새 패턴은 떠올려지지만($0.97$) 둘이 합쳐진다. 새 패턴은 남의 뉴런 $7.3$개, 옛 패턴은 $7.9$개를 끌고 나온다. $a=0.2$부터는 하나 미만 & 이 장의 유도 & 모델 \\
9 & 명제~\ref{prop:writeread}: 쓰기와 읽기가 똑같이 잘 되는 수준 $a$는 없다(그림~\ref{fig:acetsim}) & 같은 스크립트, 학습률 $\eta a$: 새 패턴은 한 번 제시로 $a\ge0.9$에서 완전히, $a=0.75$에서 $0.8$만큼 기억되고, $a\le0.5$에서는 기억되지 않는다. 절반으로부터의 읽기: $a\le0.2$에서 $1.0$, $a=0.3$에서 $0.96$, $a=0.4$에서 $0.04$ & 이 장의 유도 & 모델 \\
10 & 뇌에서는 브로드캐스트 전선 하나, 즉 \nt{ACET}가 쓰기와 읽기를 전환한다 & 이 발상은 절편 측정을 바탕으로 만든 모델에서 나왔다. 가장 직접적인 실험은 사람에서다. 무스카린성 수용체 차단제 스코폴라민은 새 단어 쌍, 특히 옛 쌍과 겹치는 쌍의 학습을 해치지만, 이미 배운 쌍을 떠올리는 것은 해치지 않는다. 네트워크의 두 모드를 직접 측정한 결과는 없다 & \cite{HasselmoAndersonBower1992,Atri2004} & 가설 \\
11 & 필터~\eqref{eq:kalman}는 최적이다. 입력 가중치와 학습률은 같은 양 $P/R$이고, 정상 상태에서 $P=\sqrt{qR}$, 기억은 $\sqrt{R/q}$이다(명제~\ref{prop:kalman}) & 실험이 필요 없다. 칼만--뷰시 필터의 최적성은 추정 이론의 고전적 결과이고, 정상 상태는 이 장에서 유도했다 & \cite{KalmanBucy1961} & 이론 \\
12 & \nt{ACET}는 $P$와 비슷한 양, 즉 예상된 불확실성을 네트워크에 알린다 & 주의의 확률 모델에서 제안되었다. 아세틸콜린 수준이 추정의 불확실성을 따라간다는 직접 측정은 없다 & \cite{YuDayan2005} & 가설 \\
13 & \nt{ACET} 뉴런의 일부는 보상과 처벌에 20밀리초 남짓 만에, 강화가 뜻밖일수록 더 강하게 응답한다 & 생쥐, 광유전학으로 식별한 기저 전뇌 콜린성 뉴런: 보상과 처벌에 대한 응답은 $18\pm3$\,ms 뒤, 크기는 강화가 뜻밖일수록 커지며, 두 핵의 뉴런에서 응답이 비슷하다 & \cite{Hangya2015} & 측정됨 \\
14 & \nt{NORA}는 세계의 예상치 못한 변화를 알아채고, \nt{ACET}는 네트워크를 쓰기 모드로 유지한다 & 이 분업은 모델에서 제안되었다. 간접 증거: 사람에서 \nt{NORA}와 연관된 동공 지름이 변화와 학습률을 따라간다. \nt{NORA}--\nt{ACET} 쌍을 직접 검증한 실험은 없다 & \cite{YuDayan2005,Nassar2012} & 가설 \\
15 & 서파 수면에서 읽기 모드의 네트워크는 낮에 기록한 것을 재생하고, 이 재생이 그것을 장기 기억으로 옮겨 쓴다 & 쥐 해마 뉴런 앙상블 기록: 낮에 함께 작동한 세포는 이어지는 서파 수면에서 더 자주 함께, 같은 순서로 발화한다(측정됨). 옮겨 쓰기에 관해서는, 학습 뒤 해마 잔물결을 억제하면 공간 기억이 나빠지고 자발적 잔물결을 늘이면 좋아진다. 그 원인이 낮은 \nt{ACET} 수준이라는 것은 증명되지 않았다 & \cite{WilsonMcNaughton1994,Skaggs1996,Girardeau2009,FernandezRuiz2019,Hasselmo1999} & 가설 \\
\bottomrule
\end{longtable}}
